% GENERATED FILE. Edit canonical lessons, then run npm run generate:books.
% canonical-source-hash: fnv1a64:57005f6854735a75
% canonical-lessons: ZH-C515-hanleng, ZH-C515-siji, ZH-W515-ji-season, ZH-C515-yuji, ZH-C515-bishu, ZH-C515-shouhuo

\chapter{Four Seasons}
\label{ch:zh-four-seasons}
\hlchaptermodality{515}

\begin{chapteropening}
\textbf{By the end of this chapter:} \emph{I can talk about the four seasons, the rainy season, bitter cold, getting away from the heat and the harvest.}

Complete the last lesson of chapter 515: I can talk about the four seasons, the rainy season, bitter cold, getting away from the heat and the harvest.
\end{chapteropening}

\section[hánlěng]{hánlěng — bitterly cold (of a climate)}

\label{lesson:ZH-C515-hanleng}



\begin{quote}
Before the new one: say the Mandarin for to be frozen stiff, then the Mandarin for to melt.
\end{quote}

\subsection*{You'll want to know: hánlěng}
\emph{hánlěng} — \textquotedblleft{}bitterly cold (of a climate)\textquotedblright{}.

\begin{quote}
\emph{zhèr de dōngtiān bǐ Běijīng hánlěng} — Winters here are colder than in Beijing.
\end{quote}

Its characters come later; for now it is a word for the ear.

\begin{cousinweb}
\emph{hánlěng} is built from \emph{hán} and \zh{冷} 'cold', both inherited from Middle Chinese, \emph{han} and \emph{lengX} (Guangyun, in Baxter's transcription). The dictionary glosses the first 'cold, as against heat'; it lists the second three times, \emph{leng}, \emph{laengX} and \emph{lengX}, the last two glossed 'cold'. Two words for cold, put together, make it bitter.
\end{cousinweb}

\begin{grammarlens}[title={What you've built}]
One more word for saying how it was.
\end{grammarlens}

\subsection*{Your turn}
\begin{itemize}
\raggedright
  \item \emph{Say it:} \emph{hánlěng}
  \item \emph{Say it:} \emph{hánlěng}, once more
  \item \emph{Say it:} \emph{zhèr de dōngtiān bǐ Běijīng hánlěng} — Winters here are colder than in Beijing.
  \item \emph{Recall:} say the Mandarin for to be frozen stiff, then the Mandarin for to melt, then say \emph{hánlěng} again
\end{itemize}

\begin{tcolorbox}[breakable,colback=teal!4,colframe=teal!35!black,title={Before you move on}]
What does hánlěng mean? (\textquotedblleft{}Bitterly cold (of a climate)\textquotedblright{}.) Say it once more.
\end{tcolorbox}
\section[\texorpdfstring{\zh{四季} (sìjì)}{四季 (sìjì)}]{\zh{四季} (sìjì) — the four seasons}

\label{lesson:ZH-C515-siji}



\begin{quote}
Before the new one: say the Mandarin for to melt, then the Mandarin for bitterly cold.
\end{quote}

\subsection*{You'll want to know: \zh{四季}}
\textbf{\zh{四季}} — \emph{sìjì} — \textquotedblleft{}the four seasons\textquotedblright{}.

\begin{quote}
\textbf{\zh{这儿四季都有花}} \emph{zhèr sìjì dōu yǒu huā} — There are flowers here all year round.
\end{quote}

The new character, in pieces: \zh{季} — eight strokes in two pieces: on top the grain piece, five strokes, as on the left of \zh{秋}, but spread wide; and \zh{子} at the foot, three strokes. The next lesson puts it in your hand.

\begin{cousinweb}
\textbf{\zh{四季}} (\emph{sìjì}) is built from \zh{四} 'four' and \zh{季}, both inherited from Middle Chinese, \emph{sijH} and \emph{kjwijH} (Guangyun, in Baxter's transcription). The dictionary glosses the second as the youngest of brothers and 'young'; how it came to mean a season, as in \zh{季节}, is not settled.
\end{cousinweb}

\begin{grammarlens}[title={What you've built}]
One more word for saying how it was.
\end{grammarlens}

\subsection*{Your turn}
\begin{itemize}
\raggedright
  \item \emph{Say it:} \emph{sìjì}
  \item \emph{Say it:} \emph{sìjì}, once more
  \item \emph{Say it:} \emph{zhèr sìjì dōu yǒu huā} — There are flowers here all year round.
  \item \emph{Recall:} say the Mandarin for to melt, then the Mandarin for bitterly cold, then say \emph{sìjì} again
\end{itemize}

\begin{tcolorbox}[breakable,colback=teal!4,colframe=teal!35!black,title={Before you move on}]
What does \zh{四季} mean? (\textquotedblleft{}The four seasons\textquotedblright{}.) Say it once more.
\end{tcolorbox}
\section[\texorpdfstring{\zh{季} (jì)}{季 (jì)}]{\zh{季} — the grain piece, then \zh{子} under it}

\label{lesson:ZH-W515-ji-season}



\begin{quote}
The word first, then the character inside it.

\begin{itemize}
\raggedright
  \item \emph{Recall:} say \emph{sìjì} — \textbf{R1}, one lesson back, before it decays
\end{itemize}

\begin{quote}
\zh{秋} \zh{子}
\end{quote}

\emph{Read it:} each one aloud, with nothing said first — \emph{qiū}, \emph{zǐ}
\end{quote}

\subsection*{Script — the grain piece, then \zh{子} under it}
\hlblockfigure{\detokenize{figures/ZH-W515-ji-season-filmstrip.pdf}}{How \zh{季} is written, stroke by stroke}

\begin{quote}
\zh{季}
\end{quote}

\textbf{Two pieces, top to bottom.} Count the pieces, not the strokes.

\textbf{The top piece, the grain piece, five strokes}, as on the left of \zh{秋}, but spread out wide: a short stroke falling to the left at the top, a bar, a short vertical down through the bar, a stroke falling to the left from the crossing, then a long stroke falling to the right.

\textbf{The bottom piece, \zh{子}, three strokes}, as you wrote it: a bar that turns down to the left, a vertical that curves down and hooks to the left, then a long bar across.

5 + 3 = eight.

\subsection*{Your turn}
\begin{itemize}
\raggedright
  \item \emph{Write it:} \zh{季}, all eight, with the model in front of you
  \item \emph{Cover:} write it again from memory, and count the pieces, not the strokes
\end{itemize}

\begin{tcolorbox}[breakable,colback=teal!4,colframe=teal!35!black,title={Before you move on}]
How many pieces does \zh{季} have? (\textbf{Two: the grain piece on top and \zh{子} at the foot}.) How many strokes? (\textbf{Eight: five and three}.)

Source: \href{https://raw.githubusercontent.com/chanind/hanzi-writer-data/68d10a4b21150cae5e1ebbd223eed289cf32d90c/data/\%E5\%AD\%A3.json}{Hanzi Writer Data \zh{季}, pinned stroke paths and medians}.
\end{tcolorbox}
\section[\texorpdfstring{\zh{雨季} (yǔjì)}{雨季 (yǔjì)}]{\zh{雨季} (yǔjì) — rainy season}

\label{lesson:ZH-C515-yuji}



\begin{quote}
Before the new one: say the Mandarin for bitterly cold, then the Mandarin for the four seasons.
\end{quote}

\subsection*{You'll want to know: \zh{雨季}}
\textbf{\zh{雨季}} — \emph{yǔjì} — \textquotedblleft{}rainy season\textquotedblright{}.

\begin{quote}
\textbf{\zh{我最不喜欢雨季}} \emph{wǒ zuì bù xǐhuan yǔjì} — I like the rainy season least of all.
\end{quote}

\begin{cousinweb}
\textbf{\zh{雨季}} (\emph{yǔjì}) is built from \zh{雨} 'rain' and \zh{季} 'a season', both inherited from Middle Chinese, \emph{hjuX} and \emph{kjwijH} (Guangyun, in Baxter's transcription); the dictionary glosses the second as the youngest of brothers, so its sense 'season' is not given there.
\end{cousinweb}

\begin{grammarlens}[title={What you've built}]
One more word for saying how it was.
\end{grammarlens}

\subsection*{Your turn}
\begin{itemize}
\raggedright
  \item \emph{Say it:} \emph{yǔjì}
  \item \emph{Say it:} \emph{yǔjì}, once more
  \item \emph{Say it:} \emph{wǒ zuì bù xǐhuan yǔjì} — I like the rainy season least of all.
  \item \emph{Recall:} say the Mandarin for bitterly cold, then the Mandarin for the four seasons, then say \emph{yǔjì} again
\end{itemize}

\begin{tcolorbox}[breakable,colback=teal!4,colframe=teal!35!black,title={Before you move on}]
What does \zh{雨季} mean? (\textquotedblleft{}Rainy season\textquotedblright{}.) Say it once more.
\end{tcolorbox}
\section[bìshǔ]{bìshǔ — to escape the summer heat}

\label{lesson:ZH-C515-bishu}



\begin{quote}
Before the new one: say the Mandarin for the four seasons, then the Mandarin for the rainy season.
\end{quote}

\subsection*{You'll want to know: bìshǔ}
\emph{bìshǔ} — \textquotedblleft{}to escape the summer heat\textquotedblright{}.

\begin{quote}
\emph{xiàtiān wǒ hé fùmǔ qù shān lǐ bìshǔ} — In summer I go up into the hills with my parents to get away from the heat.
\end{quote}

Its characters come later; for now it is a word for the ear.

\begin{cousinweb}
\emph{bìshǔ} is built from \emph{bì} 'to avoid' and \emph{shǔ} 'heat', both inherited from Middle Chinese, \emph{bjieH} and \emph{syoX} (Guangyun, in Baxter's transcription), which glosses the first 'to turn away from' and the second 'heat'.
\end{cousinweb}

\begin{grammarlens}[title={What you've built}]
One more word for saying how it was.
\end{grammarlens}

\subsection*{Your turn}
\begin{itemize}
\raggedright
  \item \emph{Say it:} \emph{bìshǔ}
  \item \emph{Say it:} \emph{bìshǔ}, once more
  \item \emph{Say it:} \emph{xiàtiān wǒ hé fùmǔ qù shān lǐ bìshǔ} — In summer I go up into the hills with my parents to get away from the heat.
  \item \emph{Recall:} say the Mandarin for the four seasons, then the Mandarin for the rainy season, then say \emph{bìshǔ} again
\end{itemize}

\begin{tcolorbox}[breakable,colback=teal!4,colframe=teal!35!black,title={Before you move on}]
What does bìshǔ mean? (\textquotedblleft{}To escape the summer heat\textquotedblright{}.) Say it once more.
\end{tcolorbox}
\section[shōuhuò]{shōuhuò — to bring in the harvest; to reap}

\label{lesson:ZH-C515-shouhuo}



\begin{quote}
Before the new one: say the Mandarin for the rainy season, then the Mandarin for to escape the summer heat.
\end{quote}

\subsection*{You'll want to know: shōuhuò}
\emph{shōuhuò} — \textquotedblleft{}to bring in the harvest; to reap\textquotedblright{}.

\begin{quote}
\emph{qiūtiān shì shōuhuò de jìjié} — Autumn is the season for bringing in the harvest.
\end{quote}

Its characters come later; for now it is a word for the ear.

\begin{cousinweb}
\emph{shōuhuò} is built from \emph{shōu} 'to gather in' and \emph{huò} 'to reap', both inherited from Middle Chinese, \emph{syuw} and \emph{hwak} (Guangyun, in Baxter's transcription), which glosses the first 'to collect' and the second, written with the grain piece, 'to cut down the crop'. The falling tone of huò is not the regular outcome of \emph{hwak}, which would give huó.
\end{cousinweb}

\begin{grammarlens}[title={What you've built}]
One more word for saying how it was.
\end{grammarlens}

\subsection*{Your turn}
\begin{itemize}
\raggedright
  \item \emph{Say it:} \emph{shōuhuò}
  \item \emph{Say it:} \emph{shōuhuò}, once more
  \item \emph{Say it:} \emph{qiūtiān shì shōuhuò de jìjié} — Autumn is the season for bringing in the harvest.
  \item \emph{Recall:} say the Mandarin for the rainy season, then the Mandarin for to escape the summer heat, then say \emph{shōuhuò} again
\end{itemize}

\begin{tcolorbox}[breakable,colback=teal!4,colframe=teal!35!black,title={Before you move on}]
What does shōuhuò mean? (\textquotedblleft{}To bring in the harvest; to reap\textquotedblright{}.) Say it once more.
\end{tcolorbox}
